% Generated draft: requires mathematical and source review.
\documentclass[11pt]{article}
\usepackage[margin=25mm]{geometry}
\usepackage{fontspec}
\setmainfont{Latin Modern Roman}
\usepackage{amsmath,amssymb,mathtools,unicode-math}
\setmathfont{Latin Modern Math}
\usepackage{enumitem,needspace,xurl,hyperref}
\setlength{\parindent}{0pt}
\setlength{\parskip}{5pt}
\setlength{\emergencystretch}{3em}
\newcommand{\N}{\mathbb{N}}
\newcommand{\Z}{\mathbb{Z}}
\newcommand{\Q}{\mathbb{Q}}
\newcommand{\R}{\mathbb{R}}
\newcommand{\C}{\mathbb{C}}
\newcommand{\F}{\mathbb{F}}
\newcommand{\A}{\mathbb{A}}
\newcommand{\PP}{\mathbb{P}}
\newcommand{\E}{\mathbb{E}}
\newcommand{\eps}{\varepsilon}
\newcommand{\dd}{\,\mathrm{d}}
\begin{document}
\section*{UnsolvedMath ID: 30004510}
\textbf{Source problem code:} catalogue entry
\subsection{Definitions and mathematical statement}
Let \(S\) be a nonempty finite state space and let
\[
\mathcal P(S)=\left\{\rho\in[0,1]^S:\sum_{i\in S}\rho_i=1\right\}
\]
be its probability simplex. A finite-state Markov chain has probability laws in \(\mathcal P(S)\). The source does not specify its time parameter. In the discrete-time interpretation, a transition matrix \(P\) satisfies \(P_{ij}\geq0\) and \(\sum_jP_{ij}=1\), and row-vector laws evolve by \(\rho_{n+1}=\rho_nP\). In the continuous-time interpretation, a generator \(Q\) satisfies \(Q_{ij}\geq0\) for \(i\ne j\) and \(\sum_jQ_{ij}=0\), and laws evolve by \(\dot\rho_t=\rho_tQ\). These are alternative interpretations, not additional hypotheses supplied by the record.

For a specified stationary distribution \(\pi\), detailed balance means \(\pi_iP_{ij}=\pi_jP_{ji}\) in discrete time, or \(\pi_iQ_{ij}=\pi_jQ_{ji}\) in continuous time, for every \(i,j\in S\). This describes the usual stationary notion of reversibility; the source does not specify the reference distribution or precisely define its term ``irreversible.''

\textbf{Source question.} Do analogous force--flux, large-deviation, and gradient-flow structures exist for finite-state Markov chains, including in the irreversible case?

\textbf{Unresolved formulation.} The source does not identify the structures to which ``analogous'' refers, define the required force--flux relation or gradient-flow representation, or specify the asymptotic family and observables for a large-deviation principle. It also leaves the time convention, structural hypotheses, and scope of the intended quantifiers unspecified. Thus this is a statement of the supplied research question; the evidence does not determine a unique, fully quantified existence assertion.
\subsection{Short English statement}
The question asks whether finite-state Markov chains, including irreversible chains, admit analogous force--flux, large-deviation, and gradient-flow descriptions. The intended analogy, compatibility requirements, and precise hypotheses are missing, so the statement requires further source clarification and human mathematical review.
\subsection{Sources}
\begin{itemize}
\item Local UnsolvedMath registry: supplied record with permanent ID \(30004510\), title \emph{Force-Flux Gradient Structures for Finite-State Markov Chains}, in the Oberwolfach Reports Open Problems collection; exact source problem code \texttt{catalogue entry}. The record gives proposed year \(2020\) and supplies no named proposer.
\item Supplied report citation: \emph{Variational Methods for Evolution (2021)}, Oberwolfach Reports, DOI \texttt{10.4171/owr/2020/29}, pp.\ 1391--1467. \url{https://doi.org/10.4171/owr/2020/29}. The background attributes the question to collected workshop open-problem discussions; no specific contribution or problem page is supplied.
\item Supplied literature reference: Mielke--Renger--Peletier, \emph{Markov Chains and Diffusion}. \url{https://arxiv.org/abs/1312.7591}.
\item Supplied literature reference: Patterson--Renger--Sharma, \emph{Journal of Statistical Physics}. \url{https://link.springer.com/article/10.1007/s10955-024-03233-8}.
\item Literature assessment checked on \texttt{2026-08-22}: the supplied curation labels the problem ``partially solved'' and states that Markov-chain large-deviation and gradient-flow structures are established for independent jump processes and several particle systems, but that no general Markov-chain theorem settles the question. This is supplied dated curation, not an independently verified literature assessment. The background expressly requests independent verification of time-sensitive status and cited claims.
\item The registry labels the statement \texttt{corrected\_verified} and reports a manual semantic, typographic, and LaTeX review asserting preservation of the target and self-containment. These are supplied review claims, not mathematical verification of this draft or resolution of its formulation gaps.
\end{itemize}
\end{document}
